\section{Introduction}
\label{sec:intro}

Production information systems are acquiring a new class of tenant. Industry
4.0 and Industry 5.0 framings place progressively more autonomy in the plant
while keeping human oversight central \cite{xu2021industry}. Beyond
the deterministic programs that have historically automated plants, large
language model (LLM) agent teams now read production telemetry, reason about
process state, and coordinate with each other over shop-floor tasks. These
agents inherit the reasoning and instruction-following capability of large
pretrained models \cite{zhao2023survey}. Recent
systems demonstrate LLM-based multi-agent orchestration on shop floors
\cite{zhao2026shopfloor}, LLM agents that operate industrial automation
systems \cite{xia2025control}, and LLM-driven control for skill-based
manufacturing \cite{koehle2026llmdriven}. Journal scopes that historically
hosted industrial informatics now explicitly name intelligent agents among
their enabling technologies. Application-oriented studies of LLM
assistants on assembly lines have appeared in this journal
\cite{llmdia2024}. The direction of travel is not in doubt: language-model
agents are moving from prototypes into the software stack that surrounds
production equipment.

Where these systems reach the equipment, they tend to encapsulate the action
inside a skill, a controller, or a validated proposal. The setpoint write
itself is different. A
supervisory setpoint write (changing a temperature setpoint, a pump speed,
a batch parameter) crosses from the information system into the physical
process. Its effect is delayed by process inertia, partially irreversible,
and intertwined with equipment-level protection that the information system
cannot see. Classical automation research has long warned about the
organizational, not merely technical, hazards of delegating control
\cite{bainbridge1983ironies}. Human-factors models place action
selection and execution at the stages where automation support changes
operator responsibility \cite{parasuraman2000model}. When the candidate
actor is an LLM agent, the delegation question becomes sharper: the agent is
capable, fallible, and susceptible to injected instructions
\cite{greshake2023injection,liu2024pinjection}.

Agent frameworks do govern actions, but at the wrong granularity for plants.
Session-level tool permissions decide \emph{whether} a tool may be called.
They do not encode \emph{plant semantics}: safe ranges, recipe windows,
batch context, and the one-proposal-per-actuator constraint are invisible
to them. Runtime
enforcement frameworks for LLM agents constrain tool calls by policies over
arguments \cite{wang2025agentspec}, and capability-based designs contain
prompt injection \cite{debenedetti2025camel}. Neither resolves which
control caught a given plant write. Neither binds a human approval to the
exact proposed operation under changing plant state, and neither
distinguishes a write that a device confirmed from one that a transport
layer merely acknowledged. On the industrial side, policy-governed agentic control
\cite{lin2026agentic} and P\&ID-grounded validation \cite{schall2026safe}
show that governance is recognized as necessary. Yet the audit question
remains open: when a write is blocked or executed, an engineer must be able
to see \emph{which layer} decided, on \emph{which configuration}, with
\emph{what receipt}.

A running example makes the unit of governance concrete. An agent observes
a melt-pressure excursion and proposes $130\to127$\,rpm. The value is
inside the node's hard engineering range, but the active recipe window and
the one-proposal-in-flight rule must also be consulted, and a second team
may already hold an opposing proposal. The platform binds the request to
the current line--product--recipe--run tuple, returns a typed verdict that
names any excluding layer, asks a human to approve the exact value and
applicable limits when the binding is manual, rechecks the tuple at
dispatch, reads the value back, and records the receipt grade and every
decision in the ledger. This is the path that Sections~\ref{sec:arch}
and~\ref{sec:eval} evaluate; the rest of the paper explains why each
boundary is necessary and what evidence supports it.

This paper takes the agent-to-plant write path itself as the unit of
governance. We present a configuration-driven platform in which LLM agent
teams and a production line share one runtime, and in which every
machine-originated setpoint write must cross a sequence of separately
inspectable boundaries. A binding model scopes agent authority to plant
nodes. A four-layer admission interval is formed by intersecting the node's
hard engineering range with parameter, product, and recipe limits. A human
approval record is bound to the exact proposed values, applicable limits,
and the active line--product--recipe--run tuple, and is rechecked at
dispatch time. A receipt grade separates readback-verified writes from
transport-acknowledged and unverified ones. Every decision on this path,
including every rejection, is recorded in an auditable ledger. The design
intentionally makes attribution and prevention separable: a layer that
cannot block still records, and a layer that blocks states why.

We evaluate the platform across five question areas, each backed by
re-runnable evidence artifacts. A five-protocol device integration matrix
(Modbus TCP, Modbus RTU, OPC UA, MQTT, HTTP) on real protocol stacks passes
21 of 21 checks, including platform-versus-device readback equality and
honest failure on disconnects. A four-arm ablation isolates which layer
provides interception of window-breaching writes: removing the soft
batch-window interlock lets six attack writes execute (all journaled), while
the hard engineering range produces zero false blocks in every arm. A
closed-loop case study on a simulator-backed OPC UA line (real protocol
stacks, simulated process physics) shows an LLM agent team (running a
9B-class open model through a standard harness) moving melt pressure
from 16.1 to 15.6\,MPa with two human approvals, an
agent-self-calibrated process gain, and an independent physical recheck. An
adversarial episode with the same weak model class ends with the platform
canceling the task after recording twelve tool-surface violations and nine
unapproved harness calls, with zero unapproved writes. Finally, the open
implementation passes 913 of 915 unit tests and reproduces a frozen
benchmark baseline bit-exactly across twelve days.

The contributions are:
\begin{enumerate}
\item \textbf{Configuration-driven binding model.} A configuration model
that binds agent authority to production context (lines, products,
recipes, runs, and nodes) and injects derived semantic context into agent
prompts, so that write scope is a configuration fact rather than a prompt
convention (Section~\ref{sec:arch-binding}).
\item \textbf{Layer-attributable governed write path.} A write path that
admits setpoint writes through a four-layer interval, bound human approval
with tuple recheck, and receipt-graded verification, and that records every
admission decision (execution or rejection) with the deciding layer
identified (Section~\ref{sec:arch-write}).
\item \textbf{Multi-level evaluation with mechanism attribution.} An
evaluation combining device-level protocol integration, a four-arm ablation
that attributes interception to a specific layer with zero false blocks, a
real LLM-agent closed-loop case study, a model-driven optimization tier, and
an adversarial episode (Section~\ref{sec:eval}).
\item \textbf{Deployable open implementation.} A single-runtime
implementation distributed as an installable package, with five industrial
connectivity protocols, fourteen agent harness adapters, 913 of 915 unit
tests passing, and bit-exact frozen-baseline reproduction
(Section~\ref{sec:impl}, Section~\ref{sec:eval-quality}).
\end{enumerate}

The remainder of this paper is organized as follows.
Section~\ref{sec:related} reviews LLM agents in industrial systems and
governance for machine-originated actions. Section~\ref{sec:problem} states
the failure model and requirements. Section~\ref{sec:arch} presents the
architecture and the governed write path. Section~\ref{sec:impl} describes
the implementation and deployment. Section~\ref{sec:eval} reports the
evaluation. Section~\ref{sec:discussion} discusses implications, limitations,
and threats to validity, and Section~\ref{sec:conclusion} concludes.
